\chapter{Conclusion} \label{chapter:conclusion}

Societal influence and personal experiences often shape the way users regard the privacy of their personal data. Regulatory frameworks like India's Digital Personal Data Protection Act (DPDPA) are formulated to shift the burden of privacy away from the users, while providing them control over the processing of their personal data. It is therefore important for the targeted population to have awareness of their rights and that this awareness is in alignment with the provisions of the Act.

In order to assess user understanding and potential pitfalls in their mental model, an exploratory survey of 259 educated, English-speaking participants was conducted. This survey aimed at primarily tackling two main provisions in the Act: blanket treatment of personal data and users' rights to have control over the processing practices that take place over their data. Results show that, in our sampled population the users displayed variances in the privacy-sensitivity of various data items based on their individual usage. This was primarily reflected in the way users treated Medical, Identity and Financial data items as extremely sensitive with near zero disagreement whereas demographic data items showed wide disagreements in terms of privacy. 

Additionally, users rated slightly below average levels of perceived control over the processing practices, with wide variation across respondents. Open-text responses indicate that this stems less from genuine moderate control than from inadequate procedural knowledge of how to exercise rights, an inability to confirm action and a trust in organizations that remains contingent rather than earned. The respondents stated that companies often use dark practices such as Roach Motel and Trick Wording to inhibit them from actually going through the already complex privacy policies, thus ensuring they remain uninformed. 

Users discussed the weak enforcement of legislation in the country that creates an accountability vacuum which is leveraged by the organizations. While it was clear that there was a slight optimism towards this new regulation, perceptions of skepticism and reservation towards the Government were primarily visible. This feeling of organizational distrust due to limited perceived control over processing practices and high privacy sensitivity explores the notion of Digital Resignation in this sample. Perceived regulatory shortcomings reinforce this notion, while solidifying the reasoning behind user frustrations. 

Overall, this study provides an exploratory basis for understanding user privacy perceptions of their personal data and its processing in the midst of a partially enforced legislation. It showcases how the Act's uniform treatment of personal data misaligns with users' contextual privacy sensitivities while doing little to improve their limited and largely unexercised perceived control arising from pre-existing organizational distrust --- a tension worth testing in a broader sample.